\section{Normalized cohort effects, aggregation and a selected donor test}
\label{s6:panel-influence}
The normalized representation in the earlier panel section has uncertainty
from both its numerators and its denominators. For a fixed cohort/date pair,
write $G$ for membership of the cohort, $C$ for eligible controls,
$D=Y_t-Y_{g-1}$, $p(x)=P(G=1\mid G+C=1,X=x)$ and
$m(x)=\E[D\mid C=1,X=x]$. Require disjoint $G,C$, no anticipation,
conditional parallel trends, overlap in the eligible population, and the
same cohort/calendar target as in \citet{callaway2021}. Other cohorts may
have $G=C=0$. Put $b=Cp/(1-p)$, $p_G=\E G$, $p_b=\E b$ and
\[
 \mu_G=\frac{\E[G(D-m)]}{p_G},\qquad
 \mu_b=\frac{\E[b(D-m)]}{p_b},\qquad \tau=\mu_G-\mu_b.
\]
At the true nuisances, $p_b=p_G$, $\mu_b=0$ and $\mu_G=\tau$.
The sample estimator normalizes each weighted mean separately. Holding
nuisances fixed, quotient differentiation yields
\begin{equation}
 \varphi_{gt}=
 \frac{G}{p_G}(D-m-\mu_G)
 -\frac{b}{p_b}(D-m-\mu_b).
 \label{s6:att-if}
\end{equation}
In particular, the treated term contains $-G\tau/p_G$. Centering an
unnormalized score after the fact does not reproduce this random-cohort
contribution.

At the true nuisances, the derivative in $m$ is
$-\E[G\,dm]/p_G+\E[b\,dm]/p_G=0$. The propensity derivative
multiplies a control residual of conditional mean zero, including the
normalizer derivative. These equalities justify the orthogonal influence
under cross-fitting and the appropriate product-rate conditions. With just
one nuisance correct, point consistency can survive, but those two
derivatives need not both vanish at the fitted limits. A doubly robust point
estimate is not automatically a doubly robust standard error; either obtain
the stronger rates/correct limits or include the first-stage influence of
the misspecified-limit procedure. Cluster the full influence by the actual
independent sampling unit. Many customers observed under one branch rollout
do not create independent branch assignments.

For an event-time average over a fixed eligible cohort set $\mathcal G$,
let $P=\sum_{g\in\mathcal G}p_g$ and
$\delta=\sum_gp_g\tau_{g,g+e}/P$. Its influence is
\begin{equation}
 \varphi_\delta=
 \sum_g\frac{p_g}{P}\varphi_{g,g+e}
 +\frac1P\sum_g G_g(\tau_{g,g+e}-\delta).
 \label{s6:aggregate-att-if}
\end{equation}
The second term accounts for estimated cohort shares. It disappears only
when those shares are fixed externally or the component effects coincide.
The first term is one sum on the same observation: controls reused across
cohorts create covariance. Forming a diagonal matrix of separately reported
standard errors loses that covariance. If eligibility itself is selected
using estimated effects, this fixed-set derivation is no longer applicable.
Financial exposure weighting adds the analogous derivative of estimated
exposures; externally fixed accounting weights use the full covariance of
the component effects without that extra term.

\subsection{What an exact synthetic-control inference route actually needs}
Synthetic-control weights interpolate pre-intervention outcomes, usually
under nonnegative weights summing to one. Their value depends on pre-fit,
factor balance, donor contamination and whether the treated unit is inside
the donor support \citep{abadie2010,abadie2021,benmichael2021,synthdid2021}.
A fitted counterfactual curve is not a confidence procedure by itself.
Consider one precise banking design: one of $J$ eligible regions is selected
uniformly for a sustained pilot, after the eligibility set and fitting
algorithm are fixed. Under the sharp null of no effect in every post-period,
the complete untreated outcome matrix is invariant to this assignment.

For each candidate treated region $j$, refit its donor weights using its
own pre-period and the remaining $J-1$ regions. Apply the same fixed
regularization and fit exclusions to every candidate. Compute a discrepancy
$T_j$, such as post-period mean squared gap divided by a pre-period error
with a fixed positive floor. With observed assignment $a$,
\[
 p_a=J^{-1}\sum_{j=1}^J\mathbf1\{T_j\ge T_a\}
\]
is super-uniform under the randomization: conditional on the fixed outcome
matrix, $a$ is uniform over the $J$ positions in its rank list. Ties make
the displayed test conservative. Serial correlation, common shocks and
weight estimation do not invalidate this finite assignment argument because
the whole paths are held fixed and the fitting algorithm is applied
symmetrically. Few eligible regions impose coarse attainable significance.

For a sharp constant additive effect $\tau$ in every post-period, subtract
$\tau$ from the actual treated path to impute the untreated matrix. For
each possible assignment add the hypothesized effect to that region, then
recompute the discrepancy centered on $\tau$, including its donor fit.
Inverting the resulting tests gives a confidence set for this \emph{sharp
constant-effect model}. It is not automatically an interval for an arbitrary
average treatment effect under heterogeneous post-period effects.

An observationally selected distressed bank is not uniformly randomized
among healthy donors. Its donor placebo ranks are then descriptive unless
another exchangeability/design argument is supplied. A residual-permutation
or Bayesian structural time-series route substitutes its own temporal or
prior/model restrictions; neither inherits the assignment proof. This is
why the product selects the above test for randomized regional pilots and
returns an explicit additional-design requirement for observational donor
studies. The method choice concerns the source of uncertainty, not a claim
that one donor-fitting algorithm is universally preferable.


A stronger temporal challenger is the null-imposed conformal construction
of \citet{conformal2022}. It imputes the hypothesized untreated outcome,
refits the counterfactual prediction under that null, and ranks its residual
discrepancy over a declared permutation group. Exact validity follows under
the required exchangeability and permutation-equivariant fitting. For
dependent observations, their Theorem 1 instead assumes a fixed postperiod
length, stationary mixing shocks with a bounded discrepancy density and
small prediction errors both on average and at the tested dates. Theorem 2
uses estimator stability plus a stationary exponentially mixing data law.
These are substantive ways to trade assignment information for temporal
prediction restrictions. Neither is a license to permute arbitrary
nonstationary bank residuals, and neither automatically changes a sharp
path hypothesis into the mean effect in Equation~\eqref{s6:aggregate-att-if}.
For an observational regional rollout this challenger can be selected only
after its null prediction and temporal restrictions are defended. The
randomized pilot route was selected here because its source of repeated
assignment is specified directly.
